\documentclass[
	10pt,								% globale Schriftgröße
	parskip=half-,						% setzt Absatzabstand hoch
	paper=a4,							% Format
	english,ngerman,					% lädt Sprachpakete
	]{scrartcl}							% Dokumentenklasse

% //////////////////// Pakete laden ////////////////////
\usepackage[fleqn]{amsmath}
\usepackage[fleqn]{mathtools}
\usepackage{amssymb}			% mathematische symbole, für \ceckmarks
\usepackage{amsthm}				% für proof
\usepackage{mathrsfs}			% für \mathscr
\usepackage{latexsym}
\usepackage{marvosym}				% für Lightning

\usepackage{fontspec} 			% funktioniert nur mit den neueren Compilern z.B. XeLaTeX
\usepackage{microtype}			% für bessere Worttrennung
\usepackage[ngerman]{babel} 	% Spracheinstellung
\usepackage{lmodern}			% verändert verwendete Schriftart, damit sie weniger pixelig ist

\usepackage{verbatim}
\usepackage{listings}			% Für Quellcode

\usepackage{graphicx}
\usepackage{tabularx}			% für Tabellen mit gleicher Spaltenbreite und automatischen Umbrüchen
\usepackage{fullpage}
\usepackage{multirow}			% für multirow in tabulars
\usepackage{rotate}
\usepackage[cmyk,table]{xcolor} % um Farben zu benutzen, kann mehr als das Paket color
\usepackage[					% Verlinkungen
	colorlinks,					% farbige Schrift, statt farbiger Rahmen
	linktocpage,				% verlinkt im Abb.Verzeichnis Seitenzahl statt Bildunterschrift
	linkcolor=blue				% setzt Farbe der Links auf blau
	]{hyperref}					% nur für digitale Anwendungen, url = "http://www.example.com"
\usepackage{url}				% für Webadressen wie e-mail usw.: "\url{http://www.example.com}"

\usepackage{enumerate}			% für versch. Aufzählungezeichen wie z.B. a)
\usepackage{xspace}				% folgt ein Leerzeichen nach einem \Befehl, wird es nicht verschluckt.
\usepackage{cancel}				% für das Durchstreichen u.a. in Matheformeln mit \cancel
\usepackage{float}              % zum Forcieren der Position von figure-Umgebungen

% zum Zeichnen (u.a. von Graphen)
\usepackage{fp}
\usepackage{tikz}
\usetikzlibrary{tikzmark}			% für \tikzmark{toRemember}
\usetikzlibrary{positioning}	% verbesserte Positionierung der Knoten
\usetikzlibrary{automata}		% für Automaten (GTI)
\usetikzlibrary{arrows}
\usetikzlibrary{shapes}
\usetikzlibrary{decorations.pathmorphing}
\usetikzlibrary{decorations.pathreplacing}
\usetikzlibrary{decorations.shapes}
\usetikzlibrary{decorations.text}

% //////////////////// Syntaxhighlighting ////////////////////
\lstloadlanguages{Python, Haskell, [LaTeX]TeX, Java}
\lstset{
   basicstyle=\footnotesize\ttfamily,	% \scriptsize the size of the fonts that are used for the code
   backgroundcolor = \color{bgcolour},	% legt Farbe der Box fest
   breakatwhitespace=false,	% sets if automatic breaks should only happen at whitespace
   breaklines=true,			% sets automatic line breaking
   captionpos=t,				% sets the caption-position to bottom, t for top
   commentstyle=\color{codeblue}\ttfamily,% comment style
   frame=single,				% adds a frame around the code
   keepspaces=true,			% keeps spaces in text, useful for keeping indentation
							% of code (possibly needs columns=flexible)
   keywordstyle=\bfseries\ttfamily\color{codepurple},% keyword style
   numbers=left,				% where to put the line-numbers;
   							% possible values are (none, left, right)
   numberstyle=\tiny\color{codegreen},	% the style that is used for the line-numbers
   numbersep=5pt,			% how far the line-numbers are from the code
   stepnumber=1,				% nummeriert nur jede i-te Zeile
   showspaces=false,			% show spaces everywhere adding particular underscores;
							% it overrides 'showstringspaces'
   showstringspaces=false,	% underline spaces within strings only
   showtabs=false,			% show tabs within strings adding particular underscores
   flexiblecolumns=false,
   tabsize=1,				% the step between two line-numbers. If 1: each line will be numbered
   stringstyle=\color{orange}\ttfamily,	% string literal style
   numberblanklines=false,				% leere Zeilen werden nicht mitnummeriert
   xleftmargin=1.2em,					% Abstand zum linken Layoutrand
   xrightmargin=0.4em,					% Abstand zum rechten Layoutrand
   aboveskip=2ex, 
}

\lstdefinestyle{py}{
   language=Python,
}
\lstdefinestyle{hs}{
   language=Haskell,
}
\lstdefinestyle{tex}{
	language=[LaTeX]TeX,
	escapeinside={\%*}{*)},     % if you want to add LaTeX within your code
	texcsstyle=*\bfseries\color{blue},% hervorhebung der tex-Schlüsselwörter
	morekeywords={*,$,\{,\},\[,\],lstinputlisting,includegraphics,
	rowcolor,columncolor,listoffigures,lstlistoflistings,
	subsection,subsubsection,textcolor,tableofcontents,colorbox,
	fcolorbox,definecolor,cellcolor,url,linktocpage,subtitle,
	subject,maketitle,usetikzlibrary,node,path,addbibresource,
	printbibliography},% if you want to add more keywords to the set
     numbers=none,
     numbersep=0pt,
     xleftmargin=0.4em,
}

\lstdefinestyle{java}{
	language=Java,
	extendedchars=true,		% lets you use non-ASCII characters;
   						% for 8-bits encodings only, does not work with UTF-8
}

\lstdefinelanguage[x64]{Assembler}     % add a "x64" dialect of Assembler
   [x86masm]{Assembler} % based on the "x86masm" dialect
   % with these extra keywords:
   {morekeywords={CDQE,CQO,CMPSQ,CMPXCHG16B,JRCXZ,LODSQ,MOVSXD, %
                  POPFQ,PUSHFQ,SCASQ,STOSQ,IRETQ,RDTSCP,SWAPGS, %
                  rax,rdx,rcx,rbx,rsi,rdi,rsp,rbp, %
                  r8,r8d,r8w,r8b,r9,r9d,r9w,r9b}
}					% for 8-bits encodings only, does not work with UTF-8

\lstdefinestyle{c}{
	language=c,
	extendedchars=true,		% for 8-bits encodings only, does not work with UTF-8
}

% //////////////////// eigene Kommandos ////////////////////
\newcommand\FU{Freie Universität Berlin\xspace}% benötigt package xspace
\newcommand\gdw{g.\,d.\,w.\xspace}
\newcommand\oBdA{o.\,B.\,d.\,A.\xspace}
\newcommand{\Eu}{\texteuro}
\newcommand\N{\mathbb{N}\xspace}
\newcommand\Q{\mathbb{Q}\xspace}
\newcommand\R{\mathbb{R}\xspace}
\newcommand\Z{\mathbb{Z}\xspace}
\newcommand\ohneNull{\ensuremath{\backslash\lbrace 0\rbrace}}% \{0}
\let\dhALT\dh	% Schreibt Befehl \dh in \dhALT um
\renewcommand\dh{d.\,h.\xspace}	%renew überschreibt command \dh
\newcommand\Bolt{\;\text{\LARGE\raisebox{-0.3em}{\Lightning}\normalsize}\xspace}% Blitz
\newcommand\zz{\ensuremath{\raisebox{+0.25ex}{Z}% zu zeigen
			\kern-0.4em\raisebox{-0.25ex}{Z}%
			\;\xspace}}
\newcommand{\from}{\ensuremath{\colon}}
\newcommand{\floor}[1]{\lfloor{#1}\rfloor}
\newcommand{\ceil}[1]{\lceil{#1}\rceil}
 \renewcommand{\L}{\ensuremath{\mathcal{L}}\xspace}
 \renewcommand{\P}{\ensuremath{\mathcal{P}}\xspace}
 \newcommand{\NL}{\ensuremath{\mathcal{N}\kern-0.2em\mathcal{L}}\xspace}
 \newcommand{\NP}{\ensuremath{\mathcal{NP}}\xspace}

% //////////////////// Mathefunktionen ////////////////////
\DeclareMathOperator{\Landau}{\mathcal{O}}
\DeclareMathOperator{\True}{True}
\DeclareMathOperator{\False}{False}

% //////////////////// eigene Theoreme ////////////////////
\newtheorem{theorem}{Satz}
\newtheorem{corollary}[theorem]{Folgerung}
\newtheorem{lemma}[theorem]{Lemma}
\newtheorem{observation}[theorem]{Beobachtung}
\newtheorem{definition}[theorem]{Definition}
\newtheorem{Literatur}[theorem]{Literatur}
% konfiguriert proof
\makeatletter
\newenvironment{Proof}[1][\proofname]{\par
  \pushQED{\qed}%
  \normalfont \topsep6\p@\@plus6\p@\relax
  \trivlist
  \item[\hskip\labelsep
%         \itshape
        \bfseries
    #1\@addpunct{.}]\ignorespaces
}{%
  \popQED\endtrivlist\@endpefalse
}
\makeatother

% //////////////////// eigene Farben ////////////////////
\let\definecolor=\xdefinecolor
\definecolor{FUgreen}{RGB}{153,204,0}
\definecolor{FUblue}{RGB}{0,51,102}

\definecolor{middlegray}{rgb}{0.5,0.5,0.5}
\definecolor{lightgray}{rgb}{0.8,0.8,0.8}
\definecolor{orange}{rgb}{0.8,0.3,0.3}
\definecolor{azur}{rgb}{0,0.7,1}
\definecolor{yac}{rgb}{0.6,0.6,0.1}
\definecolor{Pink}{rgb}{1,0,0.6}

\definecolor{bgcolour}{rgb}{0.97,0.97,0.97}
\definecolor{codegreen}{rgb}{0,0.6,0}
\definecolor{codegray}{rgb}{0.35,0.35,0.35}
\definecolor{codepurple}{rgb}{0.58,0,0.82}
\definecolor{codeblue}{rgb}{0.4,0.5,1}

% //////////////////// eigene Settings ////////////////////

\textheight = 230mm		% Höhe des Satzspiegels / Layouts
\footskip = 10ex			% Abstand zw. Fußzeile und Grundlinie letzter Textzeile
\parindent 0pt			% verhindert Einrückung der 1. Zeile eines Absatzes
\setkomafont{sectioning}{\rmfamily\bfseries}% setzt Ü-Schriften in Serifen, {disposition}											% bindet Header ein (WICHTIG)
\usepackage{graphicx}
\usepackage{amsmath}
\usepackage{amssymb}
\usepackage{fancyvrb}

\newcommand{\dozent}{Prof. R. Rojas}					% <-- Names des Dozenten eintragen
\newcommand{\projectNo}{6}
\newcommand{\veranstaltung}{Mustererkennung}
\newcommand{\semester}{WS17/18}
\newcommand{\studenten}{Boyan Hristov, Nedeltscho Petrov}
% /////////////////////// BEGIN DOKUMENT /////////////////////////


\begin{document}
% /////////////////////// BEGIN TITLEPAGE /////////////////////////
\begin{titlepage}
	\subject{\dozent}
	\title{\veranstaltung, \semester}
	\subtitle{\Large Übungsblatt \projectNo\\ \large\vspace{1ex} }
	\author{\studenten}
	\date{\normalsize \today}
\end{titlepage}

\maketitle								% Erstellt das Titelblatt
\vspace*{-9cm}							% rückt Logo an den oberen Seitenrand
\makebox[\dimexpr\textwidth+1cm][r]{	%rechtsbündig und geht rechts 1cm über Layout hinaus
	\includegraphics[width=0.4\textwidth]{src/fu_logo} % fügt FU-Logo ein
}
% /////////////////////// END TITLEPAGE /////////////////////////

\vspace{7cm}							% Abstand
\rule{\linewidth}{0.8pt}				% horizontale Linie										% erstellt die Titelseite


Link zum Git Repository: \url{https://github.com/BoyanH/FU-MachineLearning-17-18/tree/master/Solutions/Homework\projectNo}

\section*{Perceptron Classifier}


\section*{Rückgabe des Programms und damit Score und beste Annäherung}

(seed=8 in Initmethode)
\begin{lstlisting}
Smallest error for Iris-virginica vs Iris-setosa: 0.0 mistaken probes
Score Iris-virginica vs Iris-setosa: 1.0
Score Iris-virginica vs Iris-setosa on train data: 1.0
Smallest error for Iris-setosa vs Iris-versicolor: 0.0 mistaken probes
Score Iris-setosa vs Iris-versicolor: 1.0
Smallest error for Iris-virginica vs Iris-versicolor: 2.0 mistaken probes
Score Iris-virginica vs Iris-versicolor: 1.0
Score Iris-virginica vs Iris-versicolor on train data: 0.9666666666666667
Score of LDA Iris-virginica vs Iris-setosa: 1.0
Score of LDA Iris-virginica vs Iris-setosa on train data: 0.9666666666666667
\end{lstlisting}

Und mit anderem zufällig ausgewälten initialen Gewichtsvektor. (seed=1 in Initmethode)

\begin{lstlisting}
Smallest error for Iris-virginica vs Iris-setosa: 0.0 mistaken probes
Score Iris-virginica vs Iris-setosa: 1.0
Score Iris-virginica vs Iris-setosa on train data: 1.0
Smallest error for Iris-setosa vs Iris-versicolor: 0.0 mistaken probes
Score Iris-setosa vs Iris-versicolor: 1.0
Smallest error for Iris-virginica vs Iris-versicolor: 4.0 mistaken probes
Score Iris-virginica vs Iris-versicolor: 0.9
Score Iris-virginica vs Iris-versicolor on train data: 0.9555555555555556
Score of LDA Iris-virginica vs Iris-setosa: 1.0
Score of LDA Iris-virginica vs Iris-setosa on train data: 0.9666666666666667
\end{lstlisting}


\section*{Analyse}
Wie man von den Ergebnissen sieht, ist auf diesem Datensatz ein Perceptron Klassifikator fast so gut wie LDA.
Man kann aber leicht ein "multi-layer network" von Perceptrons bauen mit denen man noch bessere Ergebnisse bekommen
sollte. Weiter ist dieses Verfahren auch für "online learning" gut geigent, da man die Gewichte anhand von einzelnen
Punkten ändert. D.h. man könnte mit der Zeit von weiteren Beispielen lernen und noch bessere Ergebnisse mit der Zeit
bekommen.

Weiter sieht man, dass wir bessere Ergebnisse mit dem Testdatensatz bekommen haben, als mit dem Trainingdatensatz. Also
wir lernen nicht mehr so viel "auswendig", das Algorithmus kann besser mit sehr "untypische" Beispiele umgehen.

\section*{Implementierung}

\section*{Fit Methode}

Wir haben das vorgeschlagene Entwurf aus der Vorlesung etwas geändert. Die Implementierung ist ähnlicher zu der im
Buch "Elements of Statistical Learning". Wir nehmen die Differenz aus berechnete und erwartete Klassifizierung und
subtrahieren die von dem bisher berechnete Gewichtsvektor. Das erstetzt einfach die IF-Anweisungen, sonst gibt es keine
Unterschied. Wir haben aber gelesen, dass es besser ist, ein Lernfaktor zu benutzen und die Differenz damit zu
multiplizieren. So konvergiert man besser, da sonst ein Paar sehr "untypische" Punkte ganz schnell den Gewichtsvektor
in falsche Richtung bewegen.

Weiter werden wir maximal 500 Iterationen ausführen, bei denen nicht bessere Ergebnisse rauskommen. So haben wir auch
den Fall behandelt, in dem die Daten nicht linear separierbar sind.

Als "bessere" bzw "schlechtere" Ergebnisse haben wir am Anfang den Drehungsvinkel zwischen bisherigen und
neuen Gewichtsvektor genommen. Wir haben aber gesehen, dass nur einige "nicht typische" Punkte aus dem Datensatz
große Fehler erzeugen und man nicht bessere Ergebnisse bekommt, wenn man die korrigiert. Bei dem Datensatz war es
deutlich besser, dafür zu sorgen, dass es wenigere falsche Klassifizierungen gibt, als dass der Gewichtsvektor möglichst
nah ist an den falschen Punkten. Wir nehmen an das wird in den meisten Fällen so sein.
\begin{lstlisting}[style=py]
    def fit(self, X, y):
        t = 0
        least_error = None
        current_error = None
        best_w = self.w
        worse_iterations = 0

        while worse_iterations < 500 and (current_error is None or current_error > 0):  #  < least_error:
            if least_error is not None and current_error > least_error:
                worse_iterations += 1

            current_error = 0

            for x, yi in zip(X, y):
                error = ((self.predict_single_normalized(x) - yi)/2)
                w_new = self.w - learning_rate*error*x
                # current_error += PerceptronClassifier.get_error(w_new, self.w)
                current_error += abs(error)
                self.w = w_new

            if current_error is not None and (least_error is None or current_error < least_error):
                least_error = current_error
                best_w = np.copy(self.w)

        self.w = best_w / np.linalg.norm(best_w)
        for x, yi in zip(X, y):
            error = (self.predict_single_normalized(x) - yi)
            assert(error == 0 or least_error > 0)

\end{lstlisting}


\section*{Initialisierung}
Wir haben initial als Gewichtsvektor ein beliebigen Vektor aus dem Datensatz genommen. Alternativ sollte man auch den
Nullvektor nehmen können. Wir haben auch eine weitere Spalte in den Datensatz eingefügt mit Einsen. Das sollte im
Prinzip bessere Ergebnisse liefern, wenn es z.B. deutlich wahrscheinlicher ist, dass ein Datenpunkt zu der einen
Klasse gehört (wenn die eine Klasse häufiger vorkommt). Das sollte aber bei dem Datensatz keine große Rolle spielen.

\begin{lstlisting}[style=py]
    def __init__(self, X, y, class_a, class_b):
        ones = np.ones((len(X), 1), dtype=np.float64)
        X_normalized = np.append(ones, X, axis=1)
        y_normalized = np.vectorize(lambda x: -1 if x == class_a else 1)(y)
        np.random.seed(8)
        self.w = X_normalized[np.random.randint(0, X_normalized.shape[0], 1)][0]
        self.class_a = class_a
        self.class_b = class_b
        self.fit(X_normalized, y_normalized)
\end{lstlisting}


\section*{Vollständiges Code}


PerceptronClassifier.py
\begin{lstlisting}[style=py]
from Classifier import Classifier
import numpy as np
import math

classes_in_data_set = ['Iris-virginica', 'Iris-setosa', 'Iris-versicolor']
infinity = float('inf')
learning_rate = 0.0001

class PerceptronClassifier(Classifier):
    def __init__(self, X, y, class_a, class_b):
        ones = np.ones((len(X), 1), dtype=np.float64)
        X_normalized = np.append(ones, X, axis=1)
        y_normalized = np.vectorize(lambda x: -1 if x == class_a else 1)(y)
        np.random.seed(8)
        self.w = X_normalized[np.random.randint(0, X_normalized.shape[0], 1)][0]
        self.class_a = class_a
        self.class_b = class_b
        self.fit(X_normalized, y_normalized)

    def fit(self, X, y):
        t = 0
        least_error = None
        current_error = None
        best_w = self.w
        worse_iterations = 0

        while worse_iterations < 500 and (current_error is None or current_error > 0):  #  < least_error:
            if least_error is not None and current_error > least_error:
                worse_iterations += 1

            current_error = 0

            # what happens here is really the same as in the lecture
            # just without if statements; if e.g x is positive and we predicted negative
            # predict single would be 0, y would be 1
            # => self.w + learning_rate*x

            # learning rate is something commonly used in this algorithm, in the lecture we learned
            # a simplified version where the learning rate is 1
            for x, yi in zip(X, y):
                error = ((self.predict_single_normalized(x) - yi)/2)
                w_new = self.w - learning_rate*error*x
                # current_error += PerceptronClassifier.get_error(w_new, self.w)
                current_error += abs(error)
                self.w = w_new

            if current_error is not None and (least_error is None or current_error < least_error):
                least_error = current_error
                best_w = np.copy(self.w)

        self.w = best_w / np.linalg.norm(best_w)
        for x, yi in zip(X, y):
            error = (self.predict_single_normalized(x) - yi)
            assert(error == 0 or least_error > 0)

        print('Smallest error for {} vs {}: {} mistaken probes'.format(self.class_a, self.class_b, least_error))

    @staticmethod
    def get_error(w_new, w):
        # err in degrees rotation
        return math.acos(np.clip((w_new / np.linalg.norm(w_new)).dot(w / np.linalg.norm(w)), -1.0, 1.0))

    def project_point(self, x):
        return x.dot(self.w / np.linalg.norm(self.w))

    def predict_single_normalized(self, x):
        return 1 if self.project_point(x) > 0 else -1

    def predict_single(self, x):
        x_normalized = np.append(np.array([1]), x)
        return self.class_a if self.predict_single_normalized(x_normalized) < 0 else self.class_b

    def predict(self, X):
        return list(map(lambda x: self.predict_single(x), X))


\end{lstlisting}



PerceptronClassifierDemo.py
\begin{lstlisting}[style=py]
import numpy as np
from sklearn.discriminant_analysis import LinearDiscriminantAnalysis
from PerceptronClassifier import PerceptronClassifier, classes_in_data_set
from Parser import get_data_set, extract_classes_from_data_set

# Classes
# Iris-setosa, Iris-versicolour, Iris-virginica

X_train, X_test, y_train, y_test = get_data_set(1)

X_vi_se_train, y_vi_se_train = extract_classes_from_data_set(X_train, y_train, classes_in_data_set[:-1])
X_vi_se_test, y_vi_se_test = extract_classes_from_data_set(X_test, y_test, classes_in_data_set[:-1])
pc_vi_se = PerceptronClassifier(X_vi_se_train, y_vi_se_train, classes_in_data_set[0], classes_in_data_set[1])
score_vi_se = pc_vi_se.score(X_vi_se_test, y_vi_se_test)
score_vi_se_train = pc_vi_se.score(X_vi_se_train, y_vi_se_train)
print('Score {} vs {}: {}'.format(classes_in_data_set[0], classes_in_data_set[1], score_vi_se))
print('Score {} vs {} on train data: {}'.format(classes_in_data_set[0], classes_in_data_set[1], score_vi_se_train))


X_ve_se_train, y_ve_se_train = extract_classes_from_data_set(X_train, y_train, classes_in_data_set[1:])
X_ve_se_test, y_ve_se_test = extract_classes_from_data_set(X_test, y_test, classes_in_data_set[1:])
pc_ve_se = PerceptronClassifier(X_ve_se_train, y_ve_se_train, classes_in_data_set[1], classes_in_data_set[2])
score_ve_se = pc_ve_se.score(X_ve_se_test, y_ve_se_test)
print('Score {} vs {}: {}'.format(classes_in_data_set[1], classes_in_data_set[2], score_ve_se))

X_vi_ve_train, y_vi_ve_train = extract_classes_from_data_set(X_train, y_train, [classes_in_data_set[0],
                                                             classes_in_data_set[2]])
X_vi_ve_test, y_vi_ve_test = extract_classes_from_data_set(X_test, y_test, [classes_in_data_set[0],
                                                                            classes_in_data_set[2]])
pc_vi_ve = PerceptronClassifier(X_vi_ve_train, y_vi_ve_train, classes_in_data_set[0], classes_in_data_set[2])
score_vi_ve = pc_vi_ve.score(X_vi_ve_test, y_vi_ve_test)
score_vi_ve_train = pc_vi_ve.score(X_vi_ve_train, y_vi_ve_train)
print('Score {} vs {}: {}'.format(classes_in_data_set[0], classes_in_data_set[2], score_vi_ve))
print('Score {} vs {} on train data: {}'.format(classes_in_data_set[0], classes_in_data_set[2], score_vi_ve_train))

clf = LinearDiscriminantAnalysis()
clf.fit(X_vi_ve_train, y_vi_ve_train)
predictions = clf.predict(X_vi_ve_test)
predictions_train = clf.predict(X_vi_ve_train)
score_lda = np.mean(predictions == y_vi_ve_test)
score_lda_train = np.mean(predictions_train == y_vi_ve_train)
print('Score of LDA {} vs {}: {}'.format(classes_in_data_set[0], classes_in_data_set[1], score_lda))
print('Score of LDA {} vs {} on train data: {}'.format(classes_in_data_set[0], classes_in_data_set[1], score_lda_train))

\end{lstlisting}




Parser.py
\begin{lstlisting}[style=py]
import csv
import numpy as np
import pandas as pd
import os
from sklearn.model_selection import train_test_split


def parse_data():
    file_name = os.path.join(os.path.dirname(__file__), './Dataset/iris.data')
    return pd.read_csv(file_name, header=None).as_matrix()


def get_points_and_labels_from_data(data):
    points = np.array(data[:,:-1], dtype=np.float64)
    labels = data[:,-1]

    return points, labels


def extract_classes_from_data_set(X, y, classes):
    is_from_classes = np.vectorize(lambda y: y in classes)
    filter_arr = is_from_classes(y)
    return X[filter_arr], y[filter_arr]


def get_data_set(seed):
    data = parse_data()
    X, y = get_points_and_labels_from_data(data)
    # for determined results we use a seed for random_state, so that data is always split
    X_train, X_test, y_train, y_test = train_test_split(X, y, train_size=0.9, test_size=0.1,
                                                        random_state=seed)

    return X_train, X_test, y_train, y_test

\end{lstlisting}

Classifier.py
\begin{lstlisting}[style=py]
import numpy as np

class Classifier:
    def score(self, X, y):
        predictions = self.predict(X)
        return np.mean(predictions == y)

\end{lstlisting}


% /////////////////////// END DOKUMENT /////////////////////////
\end{document}
